\begin{abstractTH}

อุตสาหกรรมซอฟต์แวร์ที่เติบโตอย่างรวดเร็วส่งผลให้เกิดความต้องการทรัพยากรประมวลผลมหาศาล ซึ่งเป็นปัจจัยหนึ่งที่ก่อให้เกิดภาวะโลกร้อน โดยสาเหตุส่วนหนึ่งมาจากรูปแบบการเขียนโค้ดที่สิ้นเปลืองพลังงาน หรือการสร้างสิ่งที่เรียกง่า "Green Code Smells" ในโค้ด โครงงานนี้จึงพัฒนา "PyGreenSense" ซึ่งเป็นไลบรารีและส่วนขยายบน Visual Studio Code สำหรับตรวจจับ Green Code Smells ในภาษา Python เพื่อสร้างความตระหนักรู้เกี่ยวกับการเขียนโค้ดตามหลักการของวิศวกรรมซอฟต์แวร์ที่เป็นมิตรต่อสิ่งแวดล้อม (Green Software Engineering) แก่โปรแกรมเมอร์ที่สนใจผลกระทบในการเขียนโค้ดที่สิ้นเปลืองพลังงาน

ในโครงการนี้ ผู้พัฒนาได้ประยุกต์ใช้เทคนิค Abstract Syntax Tree (AST) เพื่อวิเคราะห์และระบุตำแหน่งโครงสร้างโค้ดที่สิ้นเปลืองทรัพยากร 5 ประเภท ได้แก่ God Class, Duplicated Code, Long Method, Mutable Default Arguments และ Dead Code พร้อมแสดงรายงานข้อผิดพลาดและประเมินการปล่อยคาร์บอนผ่านเครื่องมือ CodeCarbon นอกจากนี้ ยังผนวกตัวชี้วัดผลกระทบของคาร์บอนจากการกำจัด code smell เพื่อประเมินผลสัมฤทธิ์จากการปรับปรุงโค้ด ควบคู่กับการวัด Software Carbon Intensity (SCI) ต่อหน่วยฟังก์ชัน (SCI per CFP) ตามมาตรฐาน COSMIC Function Point (CFP) การทำงานเหล่านี้ช่วยให้โปรแกรมเมอร์สามารถประเมินการใช้พลังงานเทียบกับขนาดฟังก์ชันงาน และสร้างความตระหนักรู้ถึงความสัมพันธ์ระหว่างการเขียนโค้ดกับการใช้พลังงานได้อย่างเป็นรูปธรรม

\end{abstractTH}